\chapter{Implementace}
\label{ch:implementace}

Tato kapitola popisuje realizaci mobilní aplikace Foody, která navazuje na analytickou a návrhovou fázi představenou v~kapitolách~\ref{ch:analyza} a~\ref{ch:navrh}.
Na základě HiFi prototypu a funkčních požadavků definovaných v~sekci~\ref{sec:pozadavky} byla provedena implementace vybraných funkcí aplikace s~cílem ověřit, do jaké míry je navržený koncept technicky realizovatelný.
V~úvodu kapitoly je provedena rešerše přístupů k~vývoji mobilních aplikací a srovnání relevantních technologických frameworků, na jejímž základě byly zvoleny konkrétní technologie popsané v~sekci~\ref{sec:vyberTechnologii}.
Následuje popis softwarové architektury, datového modelu a integrace modelu umělé inteligence pro rozpoznávání potravin.
Dále jsou představeny vstupní modality, klíčové implementované funkce a vybrané obrazovky aplikace.
Kapitola je uzavřena přehledem stavu implementace funkčních požadavků a diskusí omezení výsledného řešení.

\section{Výběr technologií}
\label{sec:vyberTechnologii}

\subsection{Vývojová platforma}
\label{subsec:vyvojovaPlatforma}

S~ohledem na požadavek pokrýt platformy iOS i~Android při vývoji jedním autorem v~rámci diplomové práce byl zvolen multiplatformní přístup, který odstraňuje režii dvou oddělených nativních kódových základů (Swift/SwiftUI a~Kotlin/Jetpack Compose) a~riziko jejich funkční divergence.
Ze tří hlavních multiplatformních frameworků (Flutter~\cite{flutterDoc}, React Native~\cite{reactNativeDoc} a~Kotlin Multiplatform~\cite{kotlinMultiplatformDoc}) byl zvolen Flutter s~jazykem Dart, neboť nabízí nejvyšší míru sdílení kódu včetně prezentační vrstvy díky vlastnímu renderovacímu enginu zajišťujícímu konzistentní vykreslování napříč verzemi operačního systému a~disponuje zralým ekosystémem balíčků pokrývajícím všechny funkčně vyžadované oblasti (kamera, hlasový vstup, čárové kódy, lokální notifikace, integrace se zdravotními platformami, lokalizace).

Aplikace je primárně vyvíjena a~testována pro iOS s~plnou funkčností i~na Androidu díky multiplatformní povaze frameworku.
Použity jsou Flutter ve verzi 3.35 a~Dart 3.9, výchozí designový systém Material Design je doplněn o~centralizované návrhové tokeny v~souboru \texttt{lib/app\_theme.dart}.

\subsection{Databáze a perzistence dat}
\label{subsec:databaze}

Pro relační data aplikace byl zvolen framework Floor, ORM nad SQLite generující datové přístupové objekty (DAO) a databázové třídy z~anotovaných definic entit pomocí nástroje \texttt{build\_runner}~\cite{floorDoc}.
Hlavními důvody volby vůči alternativám byly podpora typově bezpečných reaktivních dotazů automaticky propagujících změny do prezentační vrstvy, podpora relačních vazeb s~kaskádovým mazáním vyžadovaná normalizovaným schématem aplikace (viz sekce~\ref{sec:datovyModel}) a jednoduchý anotační model integrující se přímočaře se vzorem DAO.
V~projektu je definováno 9~entit s~odpovídajícími 9~DAO třídami.

Pro perzistenci jednoduchých uživatelských nastavení (údaje profilu, preference, stav onboardingu) je použit balíček \texttt{shared\_preferences} zapouzdřený do dedikované služby.
API klíče pro externí služby OpenAI a Google Gemini jsou načítány ze souboru \texttt{.env} (vyloučeného z~verzovacího systému) prostřednictvím balíčku \texttt{flutter\_dotenv}.

\subsection{Model umělé inteligence a komunikace s~API}
\label{subsec:aiApi}

V~rámci návrhu aplikace Foody bylo jako klíčová technologická komponenta identifikováno propojení s~externím modelem umělé inteligence, který zajišťuje automatické rozpoznávání jídel z~fotografií, analýzu cvičení z~textových popisů a generování nutričních doporučení.
Tato podsekce popisuje výběr konkrétního poskytovatele modelu, zdůvodnění přístupu založeného na vzdáleném cloudovém API a základní komunikační vzor mezi mobilním klientem a API koncovým bodem.

Jako primární poskytovatel byl zvolen multimodální jazykový model od společnosti OpenAI, přístupný prostřednictvím rozhraní Chat Completions API.
Tento model podporuje vstup kombinující textové instrukce s~obrazovými daty kódovanými ve formátu Base64, což umožňuje předat mu fotografii jídla společně se strukturovaným promptem a kontextem uživatelského profilu v~rámci jediného API požadavku.
Odpověď modelu je strukturována ve formátu JSON definovaném v~rámci systémového promptu, přičemž obsahuje identifikaci jídla, odhad nutričních hodnot jednotlivých ingrediencí, numerickou míru nejistoty a celkový výsledek analýzy.
Pro implementaci HTTP komunikace s~API koncovým bodem je využita knihovna Dio, která poskytuje podporu pro konfiguraci hlaviček, nastavení časových limitů a strukturované zpracování chybových odpovědí~\cite{openAiApi}.

Jako alternativní poskytovatel byl do architektury aplikace integrován model Google Gemini, přístupný prostřednictvím rozhraní Generative Language API.
Architektura komunikační vrstvy je navržena tak, aby byla nezávislá na konkrétním poskytovateli, a to prostřednictvím jednotného abstraktního rozhraní pro generování odpovědí, které musí všechny implementace dodržet.
Zastřešující komponenta spravuje aktuálně zvoleného poskytovatele a umožňuje za běhu přepínat mezi implementacemi pro OpenAI a pro Gemini, přičemž obě tyto implementace splňují totéž rozhraní.
Tento návrh umožňuje rozšíření o~další poskytovatele bez zásahu do vyšších vrstev aplikace~\cite{geminiApi}.

Na straně bezpečnosti komunikace je vstup uživatele před odesláním do API ošetřen sanitizační vrstvou, která provádí sanitizaci vstupního textu (odstranění řídících znaků, ořez na maximální povolenou délku), detekci podezřelých vzorů naznačujících pokus o~\emph{prompt injection} a obalení uživatelského textu do XML delimitérů pro strukturální izolaci od systémových instrukcí (podrobněji v~sekci~\ref{subsec:promptInjection}).
Systémové prompty obsahují explicitní direktivu zakazující modelu interpretovat obsah uživatelského vstupu jako instrukce.
Tato vícevrstvá ochrana snižuje riziko manipulace chování modelu prostřednictvím záměrně formulovaných vstupů.

Přístup založený na cloudovém API s~sebou nese rovněž omezení, která je třeba zohlednit.
Aplikace vyžaduje pro funkci rozpoznávání pomocí umělé inteligence aktivní internetové připojení, přičemž zbývající funkcionality (manuální záznam, prohlížení historie, export dat) jsou plně dostupné v~režimu offline.
Latence API požadavku se v~závislosti na velikosti odesílaného obrázku a vytížení poskytovatele pohybuje řádově v~jednotkách sekund, což vyžaduje vhodný návrh uživatelského rozhraní s~indikací probíhajícího zpracování.
V~neposlední řadě je každý API požadavek zpoplatněn dle cenového modelu poskytovatele, což představuje provozní náklad, jehož výše je úměrná frekvenci používání funkcí postavených na umělé inteligenci.
Podrobný popis architektury pipeline pro umělou inteligenci, struktury promptů a mechanismu práce s~nejistotou výstupu je uveden v~sekci~\ref{sec:integraceAi}.

\subsection{Další klíčové knihovny}
\label{subsec:dalsiKnihovny}

Kromě vývojové platformy Flutter, databázové vrstvy Floor a integrace s~modely umělé inteligence využívá aplikace Foody více než 30~knihoven třetích stran z~ekosystému pub.dev~\cite{pubDev}.
Pro uživatelské funkce mají přímý dopad zejména \texttt{speech\_to\_text} pro hlasovou transkripci s~podporou češtiny i angličtiny, \texttt{mobile\_scanner} pro skenování čárových kódů, \texttt{image\_picker} pro výběr fotografií z~galerie i kamery, \texttt{flutter\_local\_notifications} pro časově plánované připomínky, \texttt{easy\_localization} pro lokalizaci uživatelského rozhraní do češtiny a angličtiny a \texttt{health} pro integraci s~Apple Health a Health Connect.

\section{Softwarová architektura}
\label{sec:softwarovaArchitektura}

Na základě technologií zvolených v~sekci~\ref{sec:vyberTechnologii} byla navržena softwarová architektura aplikace Foody optimalizovaná pro lokální zpracování dat s~komunikací s~externími službami umělé inteligence prostřednictvím REST API.
Architektura vychází z~návrhového vzoru \emph{Model-View-ViewModel} (MVVM) adaptovaného pro prostředí frameworku Flutter s~využitím reaktivních rozšíření knihovny GetX~\cite{getxDoc}.
Tato sekce popisuje celkovou architekturu systému včetně použitých architektonických vzorů a navigační model s~organizací obrazovek.
V~podsekci~\ref{subsec:celkovaArchitektura} je představena celková architektura klientské aplikace bez vlastního backendu, včetně vrstveného uspořádání komponent, použitých návrhových vzorů a struktury projektu.
V~závěrečné podsekci~\ref{subsec:navigace} je popsán navigační model aplikace, struktura hlavní obrazovky a organizace přibližně 60~obrazovek do tematických skupin.

\subsection{Celková architektura systému}
\label{subsec:celkovaArchitektura}

V~rámci návrhu architektury aplikace Foody bylo základním rozhodnutím zvolení přístupu bez vlastního backendového serveru.
Aplikace je koncipována jako samostatný mobilní klient, který veškerá uživatelská data ukládá lokálně na zařízení v~databázi SQLite a s~externími službami komunikuje přímo prostřednictvím REST API.
Tato volba reflektuje tematické zaměření práce na uživatelskou stránku aplikace a~na její experimentální ověření, nikoliv na vývoj a~provoz vlastní serverové infrastruktury.
Vlastní backend by sice aplikaci rozšířil o~možnost ukládat uživatelská data ve vzdáleném úložišti a~synchronizovat je mezi více zařízeními, z~hlediska samotné použitelnosti a~kvality výstupu pro uživatele by však neznamenal podstatnou změnu, neboť všechny klíčové výpočty probíhají buď přímo na zařízení, nebo prostřednictvím externích cloudových API.

Z~hlediska komunikace s~externími službami aplikace využívá tři vzdálená API.
Primárním poskytovatelem analýzy je rozhraní OpenAI Chat Completions API, prostřednictvím kterého jsou zpracovávány fotografické a textové vstupy pro rozpoznávání jídel, analýzu cvičení a generování nutričních doporučení.
Jako alternativní poskytovatel je integrováno rozhraní Google Gemini Generative Language API.
Tato integrace je v~kódu připravena, ale v~aktuální verzi aplikace není aktivně používána, což je v~obrázku~\ref{fig:architektura} znázorněno přerušovaným okrajem příslušného boxu.
Architektura umožňuje za běhu přepínat mezi oběma poskytovateli, jak bylo popsáno v~sekci~\ref{subsec:aiApi}.
Pro vyhledávání nutričních údajů produktů na základě čárového kódu je využíváno veřejné API služby Open Food Facts.
Veškerá síťová komunikace probíhá prostřednictvím HTTP klienta Dio, který zajišťuje jednotné zpracování časových limitů, chybových odpovědí a autorizačních hlaviček.
Schéma celkové architektury zahrnující mobilního klienta, lokální databázi a externí služby je znázorněno na obrázku~\ref{fig:architektura}.

\begin{figure}[ht]
	\centering
	\resizebox{\textwidth}{!}{% Schéma celkové architektury aplikace Foody.
% Vrstvený mobilní klient (5 vrstev) + lokální SQLite + 3 externí API.
\begin{tikzpicture}[
	font=\small\sffamily,
	node distance=4mm,
	layer/.style={
		draw,
		rounded corners=1.5pt,
		minimum height=8mm,
		minimum width=58mm,
		align=center,
		fill=black!4,
	},
	clientbox/.style={
		draw,
		thick,
		rounded corners=3pt,
		inner sep=4mm,
		fill=black!2,
	},
	ext/.style={
		draw,
		rounded corners=1.5pt,
		minimum height=11mm,
		minimum width=32mm,
		align=center,
		fill=black!6,
	},
	db/.style={
		draw,
		cylinder,
		shape border rotate=90,
		aspect=0.25,
		minimum height=14mm,
		minimum width=28mm,
		align=center,
		fill=black!6,
	},
	arr/.style={-{Latex[length=2mm]}, thick},
	dataflow/.style={font=\scriptsize\sffamily, midway, sloped, above, inner sep=1pt},
]

% Vrstvy mobilního klienta
\node[layer] (view) {Prezentační vrstva (View)\\\scriptsize \texttt{lib/screens/}, \texttt{lib/widgets/}};
\node[layer, below=of view] (ctrl) {Kontrolery (ViewModel)\\\scriptsize \texttt{lib/controller/} -- GetX};
\node[layer, below=of ctrl] (svc) {Služby a repozitáře\\\scriptsize \texttt{lib/services/}, \texttt{Repository}};
\node[layer, below=of svc] (data) {Datová vrstva\\\scriptsize \texttt{lib/database/} (Floor) -- \texttt{lib/network/} (Dio)};
\node[layer, below=of data] (model) {Modely\\\scriptsize \texttt{lib/model/} -- \texttt{json\_serializable}};

% Obal mobilního klienta -- kreslíme na background layer, aby nezakryl vnitřní vrstvy
\begin{scope}[on background layer]
\node[clientbox, fit=(view)(ctrl)(svc)(data)(model), label={[font=\small\sffamily\bfseries, anchor=south]north:Mobilní klient (Flutter)}] (client) {};
\end{scope}

% Lokální databáze pod klientem
\node[db, below=10mm of client] (sqlite) {Lokální\\databáze\\\scriptsize SQLite};

% Externí služby (vpravo) -- větší odsazení pro čitelné popisky šipek
\node[ext, right=38mm of view] (openai) {OpenAI\\\scriptsize Chat Completions API};
\node[ext, below=of openai] (gemini) {Google Gemini\\\scriptsize Generative Language API};
\node[ext, below=of gemini] (off) {Open Food Facts\\\scriptsize REST API};

% Šipky: mobilní klient ↔ SQLite (start zespodu klienta, popisek vpravo)
\draw[arr, <->] (client.south -| sqlite.north) -- (sqlite.north)
	node[pos=0.5, right=1mm, font=\scriptsize\sffamily] {SQL / DAO};

% Šipky: vrstva služeb / sítě → externí API (sběrnice ven z klienta)
\coordinate (busTop) at ($(client.east)+(6mm,0)$);
\coordinate (busOpenai) at (busTop |- openai.west);
\coordinate (busGemini) at (busTop |- gemini.west);
\coordinate (busOff)    at (busTop |- off.west);

% Popisek umístíme blíže ke sběrnici (dál od API boxu) -- pos=0.35 zleva
\draw[arr] (svc.east) -- (svc.east -| busTop) -- (busOpenai) -- node[pos=0.35, above, font=\scriptsize\sffamily] {JSON prompt + obraz} (openai.west);
\draw[arr] (svc.east) -- (svc.east -| busTop) -- (busGemini) -- node[pos=0.35, above, font=\scriptsize\sffamily] {JSON prompt} (gemini.west);
\draw[arr] (svc.east) -- (svc.east -| busTop) -- (busOff)    -- node[pos=0.35, above, font=\scriptsize\sffamily] {HTTP GET (čárový kód)} (off.west);

\end{tikzpicture}
}
	\caption{Architektura aplikace.}
	\label{fig:architektura}
\end{figure}

Vrstva kontrolerů (ViewModel) vystavuje pozorovatelné proměnné, na které se prezentační vrstva naváže a automaticky se překresluje při změně dat.
Architekturu doplňuje vzor \emph{Repository} pro abstrakci přístupu k~datům, vzor \emph{Strategy} pro zaměnitelné implementace poskytovatelů umělé inteligence (viz sekce~\ref{subsec:aiApi}) a centrální registr závislostí v~souboru \texttt{lib/locator.dart}.

V~souladu se vzorem MVVM je vnitřní architektura mobilního klienta organizována do pěti logických vrstev, které lze mapovat na tři základní role tohoto vzoru.
Nejvyšší vrstvou je prezentační vrstva (View), tvořená obrazovkami (adresář \texttt{lib/screens/}) a~sdílenými widgety (\texttt{lib/widgets/}), která definuje uživatelské rozhraní aplikace prostřednictvím deklarativní kompozice Flutter widgetů.
Pod ní se nachází vrstva kontrolerů (ViewModel, adresář \texttt{lib/controller/}), jejíž komponenty rozšiřují třídu \texttt{GetxController} a~zapouzdřují prezentační logiku jednotlivých obrazovek včetně reakcí na uživatelské akce a~správy reaktivního stavu zobrazených dat prostřednictvím pozorovatelných proměnných.

Zbývající tři vrstvy společně tvoří vrstvu Model vzoru MVVM.
První z~nich je vrstva služeb (\texttt{lib/services/}), která implementuje doménovou logiku aplikace včetně pipeline pro umělou inteligenci, správy uživatelského profilu, plánování notifikací a~sestavování doménových agregátů z~databázových entit.
Přístup k~datům zajišťují dvě oddělené vrstvy odpovídající dvěma typům datových zdrojů aplikace, přičemž abstrakci nad těmito zdroji poskytuje vzor \emph{Repository}.

Lokální datová vrstva (\texttt{lib/database/}) obsahuje definice databázových entit, datových přístupových objektů (DAO) a~migračních skriptů frameworku Floor, přičemž zajišťuje typově bezpečný přístup k~relačním datům uloženým v~SQLite.
Síťová vrstva (\texttt{lib/network/}) zapouzdřuje REST klienty pro komunikaci s~externími API a~poskytuje vyšším vrstvám typované odpovědi bez nutnosti pracovat přímo s~HTTP protokolem.
Repozitáře, jejichž ústředním zástupcem je \texttt{DayRecordRepository} pro denní záznam (sekce~\ref{subsec:repozitare}), fungují jako fasáda nad těmito datovými zdroji a~vystavují výsledky prostřednictvím reaktivních streamů, čímž kontrolery a~služby nejsou vázány na konkrétní implementaci úložiště.
Na nejnižší úrovni stojí vrstva modelů (\texttt{lib/model/}), která definuje datové třídy s~podporou serializace prostřednictvím knihovny \texttt{json\_serializable} a~slouží jako společný jazyk pro výměnu dat mezi ostatními vrstvami.

Na úrovni adresářové struktury projektu odpovídá každá architektonická vrstva samostatnému adresáři v~rámci kořenového adresáře \texttt{lib/}.
Kromě výše popsaných vrstev projekt obsahuje adresář \texttt{lib/utils/} se sdílenými pomocnými funkcemi, jako jsou definice promptů pro model umělé inteligence, validační logika a formátovací utility, a adresář \texttt{lib/generated/} s~automaticky generovanými soubory pro lokalizaci a databázový kód, které nejsou editovány ručně.
Konfigurační soubor \texttt{lib/app\_theme.dart} centralizuje veškeré návrhové tokeny aplikace (barvy, rozestupy, typografie, stíny), čímž zajišťuje konzistentní vizuální styl napříč všemi obrazovkami.
Vstupním bodem aplikace je soubor \texttt{lib/main.dart}, který řídí inicializační sekvenci zahrnující načtení konfigurace prostředí, registraci závislostí, obnovení uživatelské relace a spuštění plánovaných notifikací, přičemž samotná konfigurace \texttt{MaterialApp} je definována v~souboru \texttt{lib/app.dart}~\cite{flutterDoc}.

Komunikace mezi vrstvami probíhá přes centralizovaný registr závislostí, čímž jsou kontrolery a služby odděleny od konkrétních implementací datových zdrojů.
Tento návrh umožňuje za běhu nahradit implementaci poskytovatele umělé inteligence (sekce~\ref{subsec:aiApi}) bez zásahu do vyšších vrstev aplikace.

\subsection{Navigace a struktura obrazovek}
\label{subsec:navigace}

Pro navigaci mezi obrazovkami využívá aplikace imperativní model frameworku GetX bez pojmenované tabulky rout, což umožňuje předávat typované parametry přímo přes konstruktor cílové obrazovky~\cite{getxDoc}.
Kořenovou komponentou navigační hierarchie je hlavní obrazovka tvořená třízáložkovým shellem se záložkami Dashboard (denní přehled, seznam jídel a cvičení), Progress (týdenní a měsíční vizualizace, hmotnost) a Profile (osobní údaje, cíle, nastavení).
Přepínání záložek je řízeno reaktivní proměnnou a zachovává stav jednotlivých záložek bez nové navigace.

Plovoucí akční tlačítko ve spodní liště otevírá panel rychlých akcí s~pěti primárními vstupními body pro záznam dat (jídlo z~historie, foto, hlas, čárový kód, cvičení).
Panel je dostupný z~libovolné záložky a slouží jako centrální rozcestník pro všechny modality záznamu definované funkčními požadavky.
Aplikace obsahuje přibližně 58~obrazovek a modálních listů organizovaných v~adresáři \texttt{lib/screens/} do tematických skupin (onboarding, skenování, jídla a ingredience, logování, profil).
Modální listy jsou využívány pro doplňkové interakce nevyžadující opuštění aktuálního kontextu, jako je výběr data pro kopírování jídla, záznam hmotnosti nebo akce u~detailu cvičení.
Konkrétní navigační toky pro hlavní případy užití byly popsány v~sekci~\ref{subsec:interakcniToky} a snímky odpovídajících obrazovek jsou prezentovány v~sekci~\ref{sec:klicoveObrazovky}.

\section{Datový model}
\label{sec:datovyModel}

V~návaznosti na softwarovou architekturu představenou v~sekci~\ref{sec:softwarovaArchitektura} popisuje tato sekce konkrétní podobu datové vrstvy aplikace Foody.
Důraz je kladen na strukturu lokální relační databáze, vztahy mezi jednotlivými entitami a způsob, jakým je nad těmito entitami sestavena repozitářová vrstva poskytující doménové agregáty prezentační vrstvě.
Nejprve je v~podsekci~\ref{subsec:schemaDb} popsáno schéma databáze a vztahy mezi devíti hlavními entitami.
Podsekce~\ref{subsec:sablony} dokumentuje samostatný subsystém šablon pro opakovaně zapisované položky, který adresuje funkční požadavky na zrychlení záznamu.
Závěrečná podsekce~\ref{subsec:repozitare} popisuje repozitářovou vrstvu, transformaci databázových entit na doménové modely a mechanismus reaktivního zpřístupnění dat uživatelskému rozhraní.

\subsection{Schéma databáze a vztahy mezi entitami}
\label{subsec:schemaDb}

V~rámci datového modelu aplikace Foody je definováno celkem devět entit, které lze rozdělit do dvou logických skupin podle účelu uložených dat.
Transakční skupinu tvoří pět entit zachycujících konkrétní časově vázané záznamy uživatele (denní záznam, jídlo, ingredience, cvičení, hmotnost), zatímco kmenovou skupinu sdružují čtyři entity tvořící osobní knihovnu opakovaně používaných jídel, ingrediencí a cvičení daného uživatele, dále v~textu označované jako šablony, jejichž struktura a role jsou podrobně rozebrány v~samostatné podsekci~\ref{subsec:sablony}.
Vizuální znázornění schématu transakční skupiny s~naznačenými cizími klíči a kardinalitami je zachyceno na obrázku~\ref{fig:erDiagram}.
Kmenové entity jsou z~důvodu přehlednosti v~obrázku vynechány, neboť netvoří relační vazby s~transakční skupinou a jejich vzájemné vztahy jsou popsány v~odpovídající podsekci.

\begin{figure}[ht]
	\centering
	\resizebox{\textwidth}{!}{% ER diagram databázového schématu aplikace Foody.
% Zobrazuje pouze transakční entity (DayRecord, Meal, Ingredient, Exercise, WeightEntry).
% Šablonové (kmenové) entity jsou popsány v textu kapitoly v podsekci o systému šablon.
\begin{tikzpicture}[
	font=\scriptsize\sffamily,
	entity/.style={
		draw,
		rectangle split,
		rectangle split parts=2,
		rectangle split part fill={black!10, white},
		inner sep=3pt,
		rounded corners=1pt,
		text width=42mm,
		align=left,
		minimum height=8mm,
		outer sep=0pt,
	},
	rel/.style={-{Latex[length=2mm]}, thick},
	card/.style={
		font=\scriptsize\sffamily,
		fill=white,
		inner sep=2pt,
		outer sep=0pt,
	},
]

\node[entity] (day) {
	\textbf{DayRecord}
	\nodepart{two}
	\textit{PK} \texttt{id}\\
	\texttt{date} (UNIQUE)\\
	\texttt{calorieGoal}, \texttt{proteinGoal}\\
	\texttt{carbsGoal}, \texttt{fatGoal}
};

\node[entity, below left=18mm and -10mm of day] (meal) {
	\textbf{Meal}
	\nodepart{two}
	\textit{PK} \texttt{id}\\
	\textit{FK} \texttt{dayRecordId}\\
	\texttt{name}, \texttt{timestamp}\\
	\texttt{photoPath}, \texttt{confidence}\\
	\texttt{barcode}, \texttt{isFavorite}
};

\node[entity, below right=18mm and -10mm of day] (exercise) {
	\textbf{Exercise}
	\nodepart{two}
	\textit{PK} \texttt{id}\\
	\textit{FK} \texttt{dayRecordId}\\
	\texttt{name}, \texttt{timestamp}\\
	\texttt{durationMinutes}\\
	\texttt{caloriesBurned}, \texttt{source}
};

\node[entity, below=16mm of meal] (ingredient) {
	\textbf{Ingredient}
	\nodepart{two}
	\textit{PK} \texttt{id}\\
	\textit{FK} \texttt{mealId}\\
	\texttt{name}, \texttt{weight}, \texttt{amount}\\
	\texttt{calories}, \texttt{proteins}\\
	\texttt{carbs}, \texttt{fats}
};

\node[entity, below=16mm of exercise] (weight) {
	\textbf{WeightEntry}
	\nodepart{two}
	\textit{PK} \texttt{id}\\
	\texttt{date}, \texttt{weight}\\
	\texttt{photoPath}
};

% Vztahy -- kardinality posunuté dovnitř lajny (pos 0.25/0.75),
% aby nezasahovaly do hlaviček entit.
\draw[rel] (day.south west) -- (meal.north)
	node[card, pos=0.25] {1}
	node[card, pos=0.75] {0..N};
\draw[rel] (day.south east) -- (exercise.north)
	node[card, pos=0.25] {1}
	node[card, pos=0.75] {0..N};
\draw[rel] (meal.south) -- (ingredient.north)
	node[card, pos=0.25] {1}
	node[card, pos=0.75] {0..N};

\end{tikzpicture}
}
	\caption{ER diagram databázového schématu.}
	\label{fig:erDiagram}
\end{figure}

Transakční entity tvoří dvouúrovňovou hierarchii s~\texttt{DayRecord} v~kořeni: \texttt{Meal} je k~dennímu záznamu připojeno vztahem 1:0..N, \texttt{Ingredient} je obdobně připojena k~\texttt{Meal}, a \texttt{Exercise} tvoří paralelní podstrom přímo pod \texttt{DayRecord}.
Záznamy hmotnosti jsou modelovány jako samostatná entita bez vazby na denní záznam, což odpovídá jejich povaze nezávislého časového zápisu.
Veškeré cizí klíče jsou deklarovány s~pravidlem kaskádového mazání, které zajišťuje referenční integritu na úrovni SQLite a eliminuje riziko osiřelých záznamů při uživatelském mazání denního záznamu nebo přechodu mezi profily.

Normalizované schéma s~referenční integritou na úrovni databáze umožňuje efektivní dotazování pro analytické scénáře aplikace (týdenní a měsíční přehledy, vyhledávání oblíbených, vyhodnocování porušení dietních omezení), zatímco doménové agregáty pro uživatelské rozhraní skládá repozitářová vrstva popsaná v~podsekci~\ref{subsec:repozitare}.

\subsection{Systém šablon pro opakovaně používané položky}
\label{subsec:sablony}

V~záznamech každého uživatele se přirozeně utváří relativně stabilní okruh pravidelně konzumovaných jídel, opakovaně používaných ingrediencí a typických cvičebních aktivit.
Pro zachycení této opakovatelnosti zavádí datový model aplikace Foody samostatný subsystém šablon, jehož primárním účelem je sloužit jako osobní knihovna uživatele.
Knihovna je tvořena třemi šablonovými entitami pro jídla, ingredience a cvičení a uživatel si ji přirozeně buduje samotným používáním aplikace, bez nutnosti samostatné správy.
Na této knihovně pak stojí implementace funkčních požadavků FR-17 (oblíbené položky), FR-18 (duplikace dřívějších záznamů) a FR-19 (našeptávání názvů z~historie).

Šablona slouží jako zdroj dat při vytváření nového záznamu, aniž by s~ním byla trvale relačně vázána.
Pozdější editace šablony tak neovlivní již uložená historická data a smazání šablony z~knihovny nevyvolá ztrátu historie.
Pro řazení v~uživatelských nabídkách jsou u~každé šablony udržovány metriky frekvence použití (\texttt{usageCount}) a času posledního výskytu (\texttt{lastUsedAt}), díky kterým lze nabídnout uživateli nejprve nejčastěji a nejnověji používané položky.

Šablony jídel mají na rozdíl od ingrediencí a cvičení vnitřní hierarchii: u~každé šablony je uložen kompletní rozpis ingrediencí včetně hmotností a nutričních hodnot, takže opětovné použití obnovuje celé jídlo bez nutnosti nového volání modelu umělé inteligence nebo ručního zápisu.
Naplňování šablon probíhá automaticky při uložení nového záznamu, takže uživatel nemusí oblíbené položky spravovat samostatně.

\subsection{Repozitářová vrstva}
\label{subsec:repozitare}

Data jsou uživatelskému rozhraní zpřístupněna prostřednictvím pěti repozitářů odpovídajících věcným oblastem datového modelu (denní záznam, hmotnost, šablony jídel, ingrediencí a cvičení).
Klíčovou odpovědností \texttt{DayRecordRepository} je sestavení doménového modelu denního přehledu z~normalizovaných entit a jeho průběžná aktualizace prostřednictvím reaktivního streamu, takže denní přehled v~rozhraní se obnovuje automaticky při každé změně záznamu~\cite{getxDoc}.
Konkrétní způsob, jakým je datová vrstva využívána při integraci s~modelem umělé inteligence, je popsán v~sekci~\ref{sec:integraceAi}.

\section{Klíčové obrazovky aplikace}
\label{sec:klicoveObrazovky}

Tato sekce představuje hlavní obrazovky implementované aplikace, které tvoří uživatelskou tvář systému a ilustrují klíčové interakční toky odvozené z~případů užití v~sekci~\ref{sec:pripadyUziti}.
Pro každou obrazovku je uveden snímek z~reálného běhu aplikace na operačním systému iOS, popis její role a vazba na příslušné funkční požadavky.
Obrazovky, které ve své práci obsahují více navzájem souvisejících stavů (například hlavní obrazovka s~otevřeným kalendářem versus stav během rozpoznávání jídla), jsou prezentovány jako skupiny snímků v~jediném obrázku, aby čtenář viděl variabilitu jedné funkce vedle sebe.

\paragraph{Hlavní obrazovka.}
Hlavní obrazovka aplikace zobrazuje denní přehled kalorického příjmu (FR-02) a slouží jako rozcestník pro všechny vstupní modality záznamu.
V~horní části je umístěn scrollovatelný kalendář s~vizualizací plnění cílů pomocí barevných kruhů, pod ním souhrn kalorií vůči stanovenému cíli, vizualizace makroživin a seznam zaznamenaných jídel a cvičení.
V~dolní části je primární akční tlačítko otevírající panel rychlých akcí pro záznam jídla nebo cvičení.
Obrázek~\ref{fig:screenDashboard} zachycuje obrazovku ve třech typických stavech: ve výchozím zobrazení denního přehledu, se zobrazenou kartou rozpoznávání jídla v~průběhu volání pipeline umělé inteligence a s~otevřeným bottom sheetem kalendáře pro výběr dne.

\begin{figure}[ht]
	\centering
	\begin{minipage}[t]{0.31\textwidth}
		\centering
		\screenshot[width=\textwidth]{figures/screens/DashboardScreen_Main.png}
		\\(a)~Výchozí zobrazení
	\end{minipage}\hfill
	\begin{minipage}[t]{0.31\textwidth}
		\centering
		\screenshot[width=\textwidth]{figures/screens/DashboardScreen_RecognisingMeal.png}
		\\(b)~Stav během rozpoznávání jídla
	\end{minipage}\hfill
	\begin{minipage}[t]{0.31\textwidth}
		\centering
		\screenshot[width=\textwidth]{figures/screens/DashboardScreen_CalendarBottomSheet.png}
		\\(c)~Otevřený kalendář pro výběr dne
	\end{minipage}
	\caption{Hlavní obrazovka aplikace ve třech stavech.}
	\label{fig:screenDashboard}
\end{figure}

\paragraph{Skenování jídla.}
Obrazovka kamery slouží k~pořízení fotografie jídla pro rozpoznání umělou inteligencí (FR-06).
Rozhraní je úmyslně minimalistické, aby uživatele nerozptylovalo při focení.
Pro alternativní vstup z~galerie je k~dispozici samostatné tlačítko (FR-12).

\begin{figure}[ht]
	\centering
	\screenshot[width=0.42\textwidth]{figures/screens/ScanCameraScreen.png}
	\caption{Obrazovka skenování jídla.}
	\label{fig:screenScanCamera}
\end{figure}

\paragraph{Editace záznamu jídla.}
Editační obrazovka jídla slouží jako náhled výsledku rozpoznání pomocí umělé inteligence i jako rozhraní pro ruční úpravu nebo plně manuální vytvoření záznamu (FR-01).
Po dokončení rozpoznání obsahuje seznam navržených ingrediencí a barevnou indikaci nejistoty výstupu ve formě confidence badge (FR-08): zelená pro spolehlivost $\geq$~75~\%, žlutá pro $\geq$~50~\% a červená pro nižší hodnoty.
Uživatel může jednotlivé ingredience upravit, přidat nebo odebrat, změnit gramáž nebo počet kusů (FR-15) a v~případě potřeby vyvolat opakované rozpoznání pomocí funkce ,,Fix with AI'' (FR-13).

\begin{figure}[ht]
	\centering
	\screenshot[width=0.42\textwidth]{figures/screens/EditMealScreen.png}
	\caption{Editace záznamu jídla.}
	\label{fig:screenEditMeal}
\end{figure}

\paragraph{Výběr a záznam jídla.}
Obrazovka záznamu jídla sjednocuje vyhledávání v~historii (FR-19), oblíbené položky (FR-17), šablony jídel a duplikaci dříve uložených záznamů (FR-18) do jediného rozhraní.
Slouží jako vstupní bod pro rychlý ruční zápis i jako navigační uzel pro pokročilé varianty záznamu (skenování čárového kódu, hlasový vstup).
V~případě, že uživatel chce vytvořit zcela nové jídlo, které dosud v~historii ani ve sdílených databázích neexistuje, pokračuje z~této obrazovky na navazující obrazovku ručního vytvoření nového jídla, která je popsána detailněji v~sekci~\ref{subsec:manualniVstup}.

\begin{figure}[ht]
	\centering
	\begin{minipage}[t]{0.42\textwidth}
		\centering
		\screenshot[width=\textwidth]{figures/screens/MealLogScreen.png}
		\\(a)~Výběr z~historie a~oblíbených
	\end{minipage}\hfill
	\begin{minipage}[t]{0.42\textwidth}
		\centering
		\screenshot[width=\textwidth]{figures/screens/ManualAddNewMealScreen.png}
		\\(b)~Ruční vytvoření nového jídla
	\end{minipage}
	\caption{Manuální tok záznamu jídla.}
	\label{fig:screenMealLog}
\end{figure}

\paragraph{Hlasový vstup.}
Obrazovka hlasového vstupu umožňuje uživateli popsat jídlo přirozeným jazykem, přičemž rozpoznávání řeči probíhá v~reálném čase (FR-14).
Po dokončení nahrávání je text odeslán do pipeline pro umělou inteligenci, která vrátí strukturované rozpoznání ingrediencí.
Hlasový vstup byl ve výsledcích uživatelského testování označen jako jedna z~nejvíce oceněných funkcí aplikace.

\begin{figure}[ht]
	\centering
	\screenshot[width=0.42\textwidth]{figures/screens/VoiceLogScreen.png}
	\caption{Hlasový vstup.}
	\label{fig:screenVoiceLog}
\end{figure}

\paragraph{Sledování cvičení.}
Modul cvičení nabízí čtyři vstupní modality pro záznam aktivity: rozpoznání pomocí umělé inteligence z~textového popisu, hlasový vstup, ruční zadání a výběr ze šablon.
Obrazovka navazuje na hlavní funkcionalitu sledování příjmu energie a doplňuje ji o~výdejovou stranu energetické bilance (FR-22).

\begin{figure}[ht]
	\centering
	\screenshot[width=0.42\textwidth]{figures/screens/ExerciseLogScreen.png}
	\caption{Modul sledování cvičení.}
	\label{fig:screenExerciseLog}
\end{figure}

\paragraph{Profil a nastavení.}
Profilová obrazovka slouží jako rozcestník pro správu osobních údajů (FR-04), nutričních cílů (FR-03), dietních preferencí (FR-20), připomínek (FR-29) a integrace se zdravotními platformami.
Detailní podobrazovky jsou ilustrovány v~příslušných sekcích~\ref{sec:klicoveFunkce}.

\begin{figure}[ht]
	\centering
	\screenshot[width=0.42\textwidth]{figures/screens/ProfileScreen.png}
	\caption{Obrazovka profilu.}
	\label{fig:screenProfile}
\end{figure}

\paragraph{Dotazy v~přirozeném jazyce.}
Obrazovka dotazů v~přirozeném jazyce (Ask AI) umožňuje uživateli klást dotazy nad svými historickými nutričními daty (FR-27).
Implementace využívá dvouprůchodový systém: první volání modelu identifikuje časový rozsah dotazu, druhé pak generuje analytickou odpověď nad odpovídajícími daty z~lokální databáze.
Obrázek~\ref{fig:screenAskAi} zachycuje obrazovku ve třech typických stavech: výchozí prázdný stav, stav s~ukázkovými otázkami a stav s~vygenerovanou odpovědí.

\begin{figure}[ht]
	\centering
	\begin{minipage}[t]{0.31\textwidth}
		\centering
		\screenshot[width=\textwidth]{figures/screens/AskAIScreen.png}
		\\(a)~Výchozí stav
	\end{minipage}\hfill
	\begin{minipage}[t]{0.31\textwidth}
		\centering
		\screenshot[width=\textwidth]{figures/screens/AskAIScreenExamples.png}
		\\(b)~Ukázkové otázky
	\end{minipage}\hfill
	\begin{minipage}[t]{0.31\textwidth}
		\centering
		\screenshot[width=\textwidth]{figures/screens/AskAIScreenResponse.png}
		\\(c)~Odpověď AI
	\end{minipage}
	\caption{Dotazy v~přirozeném jazyce ve třech stavech.}
	\label{fig:screenAskAi}
\end{figure}

\paragraph{Týdenní a měsíční přehledy.}
Obrazovka přehledů vizualizuje trendy kalorického příjmu, hmotnosti a plnění cílů na týdenní a měsíční úrovni (FR-24).
Součástí je také kalendář s~vizualizací porušení dietních omezení (FR-21) a karta s~aktuální tělesnou hmotností a indexem BMI.
Vzhledem k~tomu, že obrazovka je delší než výška displeje, je v~obrázku~\ref{fig:screenProgress} prezentována ve dvou snímcích zachycujících střední a spodní část.

\begin{figure}[ht]
	\centering
	\begin{minipage}[t]{0.42\textwidth}
		\centering
		\screenshot[width=\textwidth]{figures/screens/ProgressScreen_MidPart.png}
		\\(a)~Střední část obrazovky
	\end{minipage}\hfill
	\begin{minipage}[t]{0.42\textwidth}
		\centering
		\screenshot[width=\textwidth]{figures/screens/ProgressScreen_BottomPart.png}
		\\(b)~Spodní část obrazovky
	\end{minipage}
	\caption{Obrazovka přehledů.}
	\label{fig:screenProgress}
\end{figure}

\section{Integrace modelu umělé inteligence}
\label{sec:integraceAi}

Tato sekce popisuje implementaci pipeline pro rozpoznávání jídel a cvičení pomocí externího modelu umělé inteligence, která naplňuje funkční požadavky FR-06, FR-07, FR-08, FR-10 a FR-11.
Postupně je představena celková architektura datového toku, struktura promptů a formát odpovědí, mechanismus práce s~nejistotou výstupu, vrstvená ochrana proti prompt injection a podpora více poskytovatelů modelu.
Volba multimodálního jazykového modelu místo specializované konvoluční sítě trénované na potravinovém datasetu vychází z~vývoje oboru po roce 2022, kdy multimodální modely umožnily v~jediném volání zpracovat obraz i textový popis bez nutnosti vlastního trénování~\cite{wang2026visionLanguageFood}.
Tento přístup současně obchází generalizační limit specializovaných modelů napříč různými kuchyněmi, který na transformerové architektuře CBiAFormer doložili Liu a~kol. propadem přesnosti z~92,40~\% na západním datasetu na 72,92~\% na smíšeném datasetu~\cite{liu2025deepLearning}, a~přebírá širší pokrytí trénovací distribuce dostupné u~obecných modelů~\cite{chotwanvirat2024scoping}.

\subsection{Architektura pipeline}
\label{subsec:aiPipeline}

Pipeline pro umělou inteligenci je v~aplikaci centralizována do jediné doménové služby (\texttt{AiPipelineService}), která zastřešuje celý tok od uživatelského vstupu až po finální výsledek prezentovaný v~uživatelském rozhraní.
Vstupem může být fotografie jídla doplněná o~volitelný textový popis, samostatný textový popis jídla nebo hlasově nadiktovaný popis předzpracovaný transkripční vrstvou.
Vstupní text nejprve prochází sanitizací a kontrolou bezpečnosti (viz sekce~\ref{subsec:promptInjection}) a následně je z~něj sestaven strukturovaný JSON prompt odpovídající scénáři analýzy jídla, cvičení nebo výpočtu nutričních cílů.

Sestavený prompt je odesílán přes abstraktní rozhraní, jehož konkrétní implementace odpovídá zvolenému poskytovateli modelu, a obdržená odpověď je následně parsována do typovaného doménového objektu prostřednictvím nástroje pro generování serializačního kódu.
Parsovaný výsledek prochází takzvaným \emph{confidence gate}, který na základě hodnoty spolehlivosti vrací jednu ze tří forem výsledku: úspěšné rozpoznání, rozpoznání s~varováním o~nízké jistotě, nebo selhání bez použitelného výstupu.

Vedle hlavní pipeline pro rozpoznávání jídel jsou implementovány samostatné toky pro analýzu cvičení a pro výpočet nutričních cílů.
Tok pro cvičení vychází z~textového popisu aktivity a vrací strukturu obsahující název, dobu trvání a odhadovanou energetickou hodnotu, přičemž využívá víceúrovňovou validaci popsanou v~sekci~\ref{subsec:confidence}.
Tok pro výpočet nutričních cílů aplikuje Mifflin-St~Jeorovu rovnici (rovnice~(\ref{eq:msj-muzi}) a~(\ref{eq:msj-zeny})) na uživatelovy údaje o~hmotnosti, výšce, věku, pohlaví a cíli, čímž vrací doporučené denní hodnoty kalorií a makroživin.

\begin{figure}[ht]
	\centering
	\resizebox{0.95\textwidth}{!}{% Sekvenční diagram pipeline umělé inteligence pro rozpoznání jídla z fotografie.
% Aktéři: UI -- AiPipelineService -- PromptSanitizer -- AiService -- Externí API.
\begin{tikzpicture}[
	font=\scriptsize\sffamily,
	x=1cm, y=1cm,
	actor/.style={
		draw,
		rounded corners=1.5pt,
		minimum height=8mm,
		minimum width=24mm,
		align=center,
		fill=black!8,
	},
	lifeline/.style={dashed, black!50},
	msg/.style={-{Latex[length=2mm]}, thick},
	ret/.style={-{Latex[length=2mm]}, thick, densely dashed},
	mlbl/.style={font=\scriptsize\sffamily, midway, above, sloped=false, fill=white, inner sep=1pt},
	selfnote/.style={
		draw,
		rounded corners=1pt,
		fill=black!4,
		font=\scriptsize\sffamily,
		align=left,
		inner sep=2pt,
	},
]

% --- Aktéři ---
\node[actor] (a1) at (0,0)   {Uživatelské\\rozhraní};
\node[actor] (a2) at (3.4,0) {\texttt{AiPipeline}\\\texttt{Service}};
\node[actor] (a3) at (6.8,0) {\texttt{Prompt}\\\texttt{Sanitizer}};
\node[actor] (a4) at (10.2,0){\texttt{AiService}\\(OpenAI / Gemini)};
\node[actor] (a5) at (13.4,0){Externí API\\(REST)};

% --- Lifelines ---
\foreach \a in {a1,a2,a3,a4,a5}
	\draw[lifeline] (\a.south) -- ($(\a.south)+(0,-10.5)$);

% --- Sekvence zpráv ---
% 1. UI -> Pipeline: vstup
\draw[msg] ($(a1.south)+(0,-0.5)$) -- ($(a2.south)+(0,-0.5)$)
	node[mlbl] {\texttt{analyzeFoodPhoto(photo, text)}};

% 2. Pipeline -> Sanitizer: sanitize(text)
\draw[msg] ($(a2.south)+(0,-1.4)$) -- ($(a3.south)+(0,-1.4)$)
	node[mlbl] {\texttt{sanitize(text)}};

% 3. Sanitizer -> Pipeline: cleanedText (návratová)
\draw[ret] ($(a3.south)+(0,-2.2)$) -- ($(a2.south)+(0,-2.2)$)
	node[mlbl] {\texttt{cleanedText} (regex screening, limit délky)};

% 4. Pipeline self-loop: sestavení JSON promptu
\node[selfnote, anchor=west] at ($(a2.south)+(0.3,-3.2)$) {
	\textbf{sestavení JSON promptu}\\
	(\texttt{lib/utils/prompt.dart}):\\
	\texttt{task} + \texttt{expected\_output} + \texttt{schema.rules}\\
	+ kontext uživatele (dieta, jazyk)
};
\draw[msg] ($(a2.south)+(0,-2.95)$) .. controls +(0.6,0) and +(0.6,0) .. ($(a2.south)+(0,-3.55)$);

% 5. Pipeline -> AiService: generate(prompt + image)
\draw[msg] ($(a2.south)+(0,-4.1)$) -- ($(a4.south)+(0,-4.1)$)
	node[mlbl] {\texttt{generate(prompt, imageBase64)}};

% 6. AiService -> External: HTTPS POST
\draw[msg] ($(a4.south)+(0,-4.9)$) -- ($(a5.south)+(0,-4.9)$)
	node[mlbl] {\texttt{HTTPS POST /chat/completions}};

% 7. External -> AiService: JSON response
\draw[ret] ($(a5.south)+(0,-6.4)$) -- ($(a4.south)+(0,-6.4)$)
	node[mlbl] {JSON odpověď modelu};

% 8. AiService -> Pipeline
\draw[ret] ($(a4.south)+(0,-7.2)$) -- ($(a2.south)+(0,-7.2)$)
	node[mlbl] {\texttt{AiResponse} (deserializace)};

% 9. Pipeline self-loop: confidence gate
\node[selfnote, anchor=west] at ($(a2.south)+(0.3,-8.2)$) {
	\textbf{confidence gate} (práh 0,50):\\
	\(\Rightarrow\) \texttt{success} \(\mid\) \texttt{lowConfidence} \(\mid\) \texttt{failure}
};
\draw[msg] ($(a2.south)+(0,-7.95)$) .. controls +(0.6,0) and +(0.6,0) .. ($(a2.south)+(0,-8.55)$);

% 10. Pipeline -> UI: výsledek
\draw[ret] ($(a2.south)+(0,-9.3)$) -- ($(a1.south)+(0,-9.3)$)
	node[mlbl] {\texttt{AiAnalysisResult}};

% --- Spodní rámečky aktérů (pro ukončení lifeline) ---
\foreach \a/\lbl in {a1/{},a2/{},a3/{},a4/{},a5/{}}
	\node[draw, fill=black!8, minimum height=2mm, minimum width=24mm]
		at ($(\a.south)+(0,-10.5)$) {};

\end{tikzpicture}
}
	\caption{Sekvenční diagram pipeline pro rozpoznání jídla.}
	\label{fig:aiPipeline}
\end{figure}

\subsection{Struktura promptů a formát odpovědí}
\label{subsec:prompty}

Strukturované prompty pro model umělé inteligence jsou centralizovány do jednoho modulu, který definuje šablony pro jednotlivé scénáře (rozpoznání jídla z~obrazu, rozpoznání jídla z~textu, analýza cvičení, výpočet nutričních cílů a dotazy v~přirozeném jazyce).
Plný text všech šablon je součástí zdrojového kódu odkazovaného v~příloze~\ref{app:kod}.
Každý prompt je sestaven jako JSON dokument se třemi hlavními poli: instrukcí pro model, schématem očekávané odpovědi a doplňkovými pravidly (například že numerická hodnota množství reprezentuje počet kusů a nikoli hmotnost, nebo že hodnota spolehlivosti musí ležet v~rozsahu 0~až 1).

Do promptu je zahrnut kontext uživatele, který modelu umožňuje přizpůsobit výstup individuálnímu profilu.
Mezi zahrnuté atributy patří dietní typ, vlastní dietní preference a jazykové preference, na jejichž základě model vrací názvy jídel a popisy v~odpovídajícím jazyce.
Z~uživatelského profilu se pro každý konkrétní dotaz extrahuje pouze relevantní podmnožina atributů, čímž se snižuje šum v~promptu a kontrolují se náklady na volání modelu.

Formát odpovědi pro analýzu jídla obsahuje příznak platnosti, identifikační údaje pokrmu, hodnotu spolehlivosti, údaj o~množství a strukturovaný seznam ingrediencí s~vlastními nutričními hodnotami.
Pro analýzu cvičení vrací model strukturu se jménem aktivity, dobou trvání a odhadem spálených kalorií za minutu, opět doplněnou o~hodnotu spolehlivosti.
Odpovědi jsou automaticky deserializovány do typovaných doménových objektů, čímž je práce s~výstupem na vyšších vrstvách aplikace zcela typově bezpečná.
Z~bezpečnostních důvodů popsaných v~sekci~\ref{subsec:promptInjection} je uživatelský vstup před vložením do promptu obalen do XML delimitérů.

\subsection{Práce s~nejistotou a confidence gate}
\label{subsec:confidence}

Funkční požadavek FR-08 stanovuje, že uživatel musí mít zpětnou vazbu o~spolehlivosti rozpoznání pomocí umělé inteligence.
Implementace tento požadavek naplňuje prostřednictvím takzvaného \emph{confidence gate}, který klasifikuje hodnotu spolehlivosti z~odpovědi modelu do tří úrovní a propaguje ji do uživatelského rozhraní.
Prahová hodnota pro úspěšné rozpoznání je nastavena na 0,50 jak pro analýzu jídel, tak pro analýzu cvičení.
Tato hodnota byla zvolena jako pragmatický kompromis mezi mírou falešně přijatých výstupů a frekvencí varovných hlášení, jejím podrobnějším laděním vůči empirickým datům se tato práce blíže nezabývala.

V~uživatelském rozhraní je nejistota vizualizována tříúrovňovou škálou v~rámci barevného indikátoru spolehlivosti na editační obrazovce jídla (viz obrázek~\ref{fig:screenEditMeal}).
Hodnoty $\geq$~75~\% jsou zobrazeny zeleně, hodnoty v~rozsahu 50~\% až 75~\% žlutě a hodnoty pod 50~\% červeně.
Tato barevná konvence prostupuje celou aplikaci.
V~uživatelském testování ji všichni participanti pochopili bez explicitního vysvětlování (viz sekce~\ref{sec:vysledkyTestovani}).

Logika rozhodování v~\emph{confidence gate} rozlišuje tři výsledné stavy.
Stav \emph{úspěch} odpovídá hodnotě spolehlivosti nad prahem a značí, že výsledek lze prezentovat uživateli bez zvláštního varování.
Stav \emph{nízká jistota} odpovídá hodnotě pod prahem 0,50, ale s~platným strukturovaným výstupem.
Data jsou v~tomto případě zobrazena s~viditelným varováním a doporučením kontroly.
Stav \emph{selhání} reprezentuje situaci, kdy model vrátí neplatný, prázdný nebo neparsovatelný výsledek, a aplikace v~tomto případě nabízí uživateli textový fallback (FR-11) nebo opakování pokusu.
Pro analýzu cvičení je validace doplněna o~kontrolu prázdného názvu, čímž je odlišeno tvrdé selhání od pouhé nízké jistoty.

\subsection{Ochrana proti prompt injection}
\label{subsec:promptInjection}

Vstup uživatele do pipeline pro umělou inteligenci představuje potenciální vektor útoku známý jako \emph{prompt injection}, kdy záměrně formulovaný vstup může změnit nebo přepsat původní systémové instrukce zaslané modelu~\cite{greshake2023indirect, owasp2025llm01}.
Hlavní motivací ochrany v~kontextu této aplikace není utajení samotného systémového promptu, jehož text je veřejně dostupný v~repozitáři aplikace, ale zachování předvídatelnosti chování modelu, na niž je vázána důvěra uživatele v~zobrazené nutriční odhady.
Z~pohledu uživatele lze identifikovat tři kategorie hrozeb, kterými jsou vědomý explicitní útok formulovaný jako adversariální instrukce, falešný poplach vyvolaný legitimním popisem jídla nebo cvičení s~náhodnou podobností pokusu o~injekci a obfuskovaný útok cíleně navržený tak, aby unikl jednoduché detekci.

Implementace adresuje tyto kategorie vícevrstvou ochranou v~rámci dedikovaného sanitizačního modulu (\texttt{PromptSanitizer}) volaného před každým odesláním vstupu do pipeline.
První vrstva (L1) provádí deterministické mechanické úpravy textu, kterými jsou odstranění bílých znaků na okrajích, odstranění řídících znaků a literálních značek vstupního delimitéru a oříznutí textu na maximální povolenou délku 500~znaků.
Druhá vrstva (L2) klasifikuje vstup pomocí šesti regulárních výrazů a vrací jeden ze tří stavů, kterými jsou \emph{čistý}, \emph{ambivalentní} a \emph{explicitní útok}.
Čtyři vzory zachycují fráze jednoznačně adversariální (například \uv{ignore previous instructions}, \uv{disregard your instructions}, \uv{reveal the system prompt} nebo \uv{do not respond in json}) a vstup s~takovým výskytem je odmítnut okamžitě.
Zbývající dva vzory označují \emph{ambivalentní} formulace jako \uv{you are}, \uv{act as} nebo zmínky o~systémovém promptu, u~nichž nelze vyloučit benigní výskyt v~přepisu hlasového vstupu nebo v~meta-dotazu na rozhraní Ask AI.
Vstup, který nezachytí žádný ze vzorů, je klasifikován jako \emph{čistý} a postupuje dále bez dodatečných kontrol.

Pro ambivalentní vstupy je aktivována třetí vrstva (L3), takzvaný \emph{LLM-based pre-screening}, který předá text menšímu a rychlejšímu modelu (GPT-5.4-mini) s~jediným úkolem rozhodnout, zda se jedná o~legitimní popis nebo o~pokus o~injekci.
Pokud LLM klasifikuje vstup jako útok, aplikace jej odmítne stejným způsobem jako u~explicitního vzoru.
V~opačném případě vstup pokračuje do hlavního modelu.
Nezávisle na výsledku klasifikace prochází každý legitimní vstup závěrečnou transformací, při níž je obalen do XML delimitéru \texttt{<user\_input>} pro strukturální izolaci od systémových instrukcí a do systémového promptu je doplněna explicitní anti-injekční direktiva zakazující modelu interpretovat obsah uživatelského bloku jako instrukce.
Tato direktiva tvoří průřezovou ochranu uplatňovanou vždy a slouží jako poslední linie obrany pro případ, že by adversariální vstup obě předchozí vrstvy obešel.
Vícevrstvý přístup vychází z~doporučení OWASP pro aplikace postavené na velkých jazykových modelech a z~empirických zjištění Liu et al.~\cite{owasp2025llm01, liu2024formalizing}, podle nichž žádná jednotlivá vrstva neposkytuje úplnou ochranu, ale jejich kombinace významně snižuje pravděpodobnost úspěšného útoku.

Specifickým rozhodnutím L3 vrstvy je její záměrný režim \emph{fail-open}, v~rámci kterého je vstup při výpadku síťové infrastruktury nebo nedostupnosti API předzvolovacího modelu považován automaticky za legitimní.
Tato volba reprezentuje vědomý kompromis mezi přísností bezpečnostní kontroly a dostupností služby pro uživatele, neboť přechodný technický problém na straně předzvolovacího modelu by jinak blokoval i~zcela neškodné popisy jídla a vedl by k~frustraci, která je v~kontextu každodenního používání aplikace nepřijatelná.
Reziduální riziko vyplývající z~tohoto kompromisu je sníženo vícevrstvou architekturou popsanou výše, ve které samostatně působí již vrstvy L1 a L2 a globálně XML delimiter s~anti-injekční direktivou.

Detekce založená na regulárních výrazech má z~principu nenulovou míru jak falešně pozitivních, tak falešně negativních klasifikací a cílený útočník dokáže jednoduché vzory obejít obfuskací, překlepy nebo kódováním v~Unicode~\cite{liu2024formalizing}.
Maximálním důsledkem úspěšného obejití všech vrstev je manipulace s~výstupem jediné aktuální analýzy, tedy například nesprávný odhad kalorické hodnoty nebo zamítavá odpověď modelu, neboť aplikace neukládá žádný serverový stav ani perzistentní datový kontext, který by útok mohl ovlivnit napříč relacemi.
Únik systémového promptu nepředstavuje samostatnou hrozbu, jelikož jeho text je veřejně dostupný v~repozitáři aplikace.
Pravděpodobnost falešně pozitivních klasifikací u~reálných uživatelských popisů jídla a cvičení v~přirozeném jazyce je naopak očekávána jako blízká nule, neboť tyto popisy fráze odpovídající explicitním adversariálním vzorům přirozeně neobsahují.
Vlastním podrobnějším vyhodnocením odolnosti popsaného schématu se tato práce blíže nezabývala.

Schematické zobrazení vrstev L1--L3, dvou rozhodovacích bodů a~průřezové fáze \emph{prompt assembly} je uvedeno na obrázku~\ref{fig:promptInjection}.

\begin{figure}[!ht]
	\centering
	\resizebox{0.7\textwidth}{!}{% Stavový diagram průchodu uživatelského vstupu vícevrstvou ochranou proti prompt injection.
% L1 = deterministická sanitizace (mechanika), L2 = klasifikace vzorů (regex, dvouúrovňový výstup),
% L3 = LLM pre-screening (pouze pro ambivalentní vstupy). Anti-injection direktiva + XML wrap
% jsou průřezová transformace aplikovaná na každý legitimní vstup.
\begin{tikzpicture}[
	font=\small\sffamily,
	node distance=7mm,
	state/.style={
		draw,
		rounded corners=1.5pt,
		minimum height=11mm,
		minimum width=58mm,
		align=center,
		fill=black!4,
	},
	process/.style={
		draw,
		rounded corners=1.5pt,
		minimum height=12mm,
		minimum width=66mm,
		align=center,
		fill=black!6,
	},
	assembly/.style={
		draw,
		dashed,
		rounded corners=1.5pt,
		minimum height=12mm,
		minimum width=66mm,
		align=center,
		fill=blue!4,
	},
	decision/.style={
		draw,
		diamond,
		aspect=2.2,
		minimum height=10mm,
		minimum width=34mm,
		align=center,
		fill=black!6,
		inner sep=1pt,
	},
	rejectbox/.style={state, fill=red!10, draw=red!50!black, minimum width=30mm},
	successbox/.style={state, fill=green!12, draw=green!50!black},
	arr/.style={-{Latex[length=2mm]}, thick},
	lbl/.style={font=\scriptsize\sffamily, midway, fill=white, inner sep=1pt},
]

% --- Hlavní vertikální řetězec ---
\node[state] (input) {Uživatelský vstup\\\scriptsize text / přepis hlasu / popis fotografie};

\node[process, below=of input] (l1) {%
	\textbf{L1: deterministická sanitizace}\\
	{\scriptsize ořez délky, odstranění řídících znaků, strip delimitérů}%
};

\node[process, below=of l1] (l2) {%
	\textbf{L2: klasifikace vstupních vzorů}\\
	{\scriptsize regex screening, tříúrovňový výstup}%
};

\node[decision, below=10mm of l2] (d1) {Explicitní\\útočný vzor?};

\node[decision, below=12mm of d1] (d2) {Ambivalentní\\vzor?};

\node[process, below=12mm of d2] (l3) {%
	\textbf{L3: LLM pre-screening}\\
	{\scriptsize klasifikace menším modelem (GPT-5.4-mini)}%
};

\node[decision, below=of l3] (d3) {Legitimní\\popis?};

\node[assembly, below=12mm of d3] (assembly) {%
	\textbf{Prompt assembly}\\
	{\scriptsize obalení v~XML delimitérech + anti-injection direktiva}%
};

\node[successbox, below=of assembly] (model) {Hlavní AI model\\\scriptsize zpracování analytického dotazu};

% --- Terminální stavy "Odmítnuto" (vlevo) ---
\node[rejectbox, left=18mm of d1] (rej1) {Odmítnuto};
\node[rejectbox, left=18mm of d3] (rej3) {Odmítnuto};

% --- Šipky hlavního toku ---
\draw[arr] (input) -- (l1);
\draw[arr] (l1) -- (l2);
\draw[arr] (l2) -- (d1);
\draw[arr] (d1.south) -- node[lbl] {ne} (d2.north);
\draw[arr] (d2.south) -- node[lbl] {ano} (l3.north);
\draw[arr] (l3.south) -- (d3.north);
\draw[arr] (d3.south) -- node[lbl] {ano} (assembly.north);
\draw[arr] (assembly.south) -- (model.north);

% --- Reject šipky ---
\draw[arr] (d1.west) -- node[lbl] {ano} (rej1.east);
\draw[arr] (d3.west) -- node[lbl] {ne} (rej3.east);

% --- Bypass L3: ne z D2 přeskočí přímo na prompt assembly (vpravo) ---
\coordinate (bypassTop) at ($(d2.east)+(22mm,0)$);
\coordinate (bypassBot) at (bypassTop |- assembly.east);
\draw[arr] (d2.east) -- (bypassTop) -- (bypassBot) -- (assembly.east);
\node[font=\scriptsize\sffamily, fill=white, inner sep=1pt, above]
	at ($(d2.east)!0.5!(bypassTop)$) {ne};

\end{tikzpicture}
}
	\caption{Stavový diagram ochrany proti prompt injection.}
	\label{fig:promptInjection}
\end{figure}

\subsection{Volba poskytovatele a modelu}
\label{subsec:vicePoskytovatelu}

Jako poskytovatel multimodálního jazykového modelu byla zvolena společnost OpenAI, která v~aktuálním kontextu představuje jednu z~nejdostupnějších a~ekosystémově nejstabilnějších platforem se širokou komunitní podporou a~ověřeným strukturovaným JSON výstupem.
Vzhledem k~tomu, že tato práce je tematicky zaměřena na HCI a~nikoliv na vývoj a~trénování modelů umělé inteligence, byl preferován obecný univerzální model před modelem doménově specializovaným.
Lze přepokládat, že model dotrénovaný výhradně na potravinových datech by pro doménu rozpoznávání jídel poskytoval přesnější odhady.
Jeho výměna je však v~rámci zvolené architektury proveditelná bez zásahu do vyšších vrstev aplikace a~představuje přirozený krok před případným nasazením do produkce.

Komunikační vrstva pro analýzu jídel je proto navržena tak, aby byla nezávislá na konkrétním poskytovateli, a to prostřednictvím jednotného abstraktního rozhraní pro generování odpovědí.
V~projektu jsou pro hlavní pipeline připraveny dvě strategie, a~to pro OpenAI Chat Completions a~pro Google Gemini Generative Language.
Zbývající scénáře (analýza cvičení, výpočet nutričních cílů, rozhraní Ask AI) jsou v~aktuální verzi vázány přímo na OpenAI a~rozšíření jejich abstrakce je formulováno jako námět pro další iteraci.

Empirické porovnání tří variant modelů OpenAI (GPT-5.4, GPT-5.4-mini, GPT-5.5) ve dvou iteracích systémového promptu bylo provedeno v~rámci kvantitativního benchmarku popsaného v~sekci~\ref{subsec:benchmarkPresnost}, na jehož základě byl pro produkční nasazení zvolen jako výchozí model GPT-5.4 s~promptem improved v2.

\section{Vstupní modality pro záznam stravy}
\label{sec:vstupniModality}

V~návaznosti na výsledky uživatelské studie (sekce~\ref{sec:uzivatelskaStudie}), která identifikovala potřebu rychlého a flexibilního zápisu jídel v~různých kontextech, implementuje aplikace celkem šest vstupních modalit pro záznam stravy a~doplňkově dvě modality pro záznam cvičení.
Jednotlivé modality jsou volány z~panelu rychlých akcí na hlavní obrazovce a~pokrývají funkční požadavky FR-01, FR-06, FR-11, FR-12, FR-14 a~FR-16.
Modality lze rozdělit do dvou skupin podle způsobu zpracování vstupu: AI-asistované cesty (foto, hlas, text, OCR nutriční etikety a~hlasový vstup pro cvičení) procházejí sdílenou pipeline pro umělou inteligenci a~následně editační obrazovkou, zatímco přímé cesty (čárový kód s~vyhledáním v~databázi Open Food Facts a~manuální zápis jídla i~cvičení) editační obrazovku obcházejí nebo využívají v~upravené podobě.
Tato sekce postupně popisuje implementaci každé z~nich, od plně automatizovaného fotografického vstupu přes hlasový a~textový vstup s~podporou umělé inteligence až po manuální zadání bez asistence umělé inteligence.

\subsection{AI-asistovaný vstup: foto, hlas a text}
\label{subsec:aiVstup}

Trojice modalit s~asistencí umělé inteligence sdílí společnou infrastrukturu: vstup v~podobě obrazu nebo textu je předán do pipeline pro umělou inteligenci a výsledek rozpoznání je prezentován na editační obrazovce jídla (viz obrázek~\ref{fig:screenEditMeal}) zároveň sloužící jako náhled i jako rozhraní pro úpravy.
Modality se liší pouze typem vstupu a uživatelským kontextem, ve kterém jsou volány.

Fotografický vstup (FR-06, FR-12) představuje hlavní inovační prvek aplikace a podporuje dva zdroje obrazu: živý vstup z~kamery prostřednictvím obrazovky skenování a import existujících fotografií z~galerie pomocí knihovny \texttt{image\_picker}.
Uživatel může k~fotografii volitelně přidat krátký textový popis (například \uv{kuřecí salát s~rýží}), který je odeslán společně s~obrazem.
Při prvním použití je zobrazena pětistránková onboardingová sekvence (FR-09) vysvětlující limity rozpoznávání a tipy pro pořízení kvalitního snímku (osvětlení, úhel pohledu, oddělení složek na talíři).
Tato instruktáž zároveň přispívá k~rozlišení mezi limitem modelu a technickou chybou aplikace (FR-10).

Hlasový vstup (FR-14) je realizován samostatnou transkripční službou postavenou nad knihovnou \texttt{speech\_to\_text} s~podporou češtiny a angličtiny podle aktuální lokalizace.
Na obrazovce hlasového vstupu (viz obrázek~\ref{fig:screenVoiceLog}) uživatel zahájí nahrávání stiskem tlačítka, vidí transkripci v~reálném čase a po dokončení ji potvrdí nebo upraví před odesláním do pipeline.
Tok pokrývá záznam jídla i analogický scénář pro záznam cvičení a v~uživatelském testování byl označen za jednu z~nejvíce oceněných funkcí aplikace (viz sekce~\ref{sec:doporuceni}).

Textový vstup (FR-14, FR-11) slouží jako přirozený záložní kanál pro situace, kdy fotografický vstup není praktický nebo selže (jídlo již zkonzumované, omezené podmínky pro focení).
Uživatel popíše jídlo libovolným způsobem v~přirozeném jazyce (například \uv{kuřecí prsa s~rýží a brokolicí, asi 250~gramů}) a vstup prochází stejnou pipeline jako fotografický nebo hlasový.

\subsection{Manuální zadání bez umělé inteligence}
\label{subsec:manualniVstup}

Manuální vstup naplňuje funkční požadavek FR-01 a je určen pro situace, kdy uživatel přesně zná složení jídla a nepotřebuje asistenci umělé inteligence (například domácí vaření s~váhou, balené potraviny s~nutriční etiketou nebo offline režim, kdy není dostupné připojení k~poskytovateli umělé inteligence).
Tok začíná na obrazovce záznamu jídla, která slouží jako sjednocený rozcestník pro výběr existujícího jídla z~historie nebo vytvoření nového záznamu.
V~případě nového záznamu pokračuje uživatel na obrazovku ručního vytvoření nového jídla (viz obrázek~\ref{fig:screenMealLog}), kde vyplní název jídla a postupně přidává jednotlivé ingredience.

Pro každou ingredienci uživatel zadává nutriční hodnoty (kalorie, bílkoviny, sacharidy, tuky) a hmotnost v~gramech nebo počet kusů (FR-15).
Implementace podporuje předvolby porcí (například 100~g, 1~ks, vlastní jednotky) a zobrazování zlomkových hodnot (například 1/2 porce), čímž zjednodušuje zápis běžných situací.
Detail editace ingrediencí probíhá na samostatných obrazovkách editace a detailu ingredience, které sdílí komponenty s~rozhraním pro úpravu výsledku rozpoznání umělou inteligencí.

\subsection{Čtečka čárových kódů}
\label{subsec:carovyKod}

Čtečka čárových kódů naplňuje funkční požadavek FR-16 a je realizována pomocí knihovny \texttt{mobile\_scanner} podporující formáty EAN-8, EAN-13, UPC-A, UPC-E a Code-128.
Po naskenování kódu volá vyhledávací služba klienta veřejné databáze Open Food Facts, který vrací strukturovaný výsledek obsahující název produktu, výrobce a kompletní nutriční hodnoty na 100~g.
Koordinaci skenování, vyhledávání a zobrazení výsledku zajišťuje samostatný kontroler.

V~návaznosti na funkční požadavek FR-10 (rozlišení limitu umělé inteligence a technické chyby aplikace) je v~rámci čárových kódů implementováno šest typů chyb pokrývajících situace jako neexistující kód, produkt nenalezen v~databázi, nepřipojení k~síti, časový limit dotazu, neplatný formát kódu a poškození čtečky.
Každý typ chyby je uživateli prezentován s~konkrétním vysvětlením a doporučením dalšího kroku.
Hlavním omezením této modality je závislost na kompletnosti databáze Open Food Facts, kterou v~uživatelském testování dva ze tří participantů identifikovali jako problematickou pro některé lokální české produkty (viz sekce~\ref{subsec:uiProblemyTesting}).

\subsection{Skenování nutričních etiket}
\label{subsec:nutricniEtikety}

Skenování nutričních etiket představuje doplňkovou modalitu zaměřenou na situace, kdy uživatel má v~ruce balený produkt s~nutriční tabulkou, ale jeho čárový kód není dostupný v~databázi Open Food Facts.
Implementace využívá multimodální schopnosti hlavního modelu umělé inteligence prostřednictvím specializovaného režimu skenování etiket, kdy je obrazovka kamery přepnuta do tohoto režimu a získaný obraz je odeslán do pipeline s~upraveným promptem zaměřeným na extrakci nutričních hodnot z~tabulky.

Aktuální stav implementace pokrývá základní extrakci hodnot pro 100~g produktu (kalorie, bílkoviny, sacharidy, tuky), zatímco rozšířené atributy jako vláknina, sůl nebo vybrané mikronutrienty jsou ponechány jako námět pro budoucí iteraci.
Hlavním omezením modality je rozmanitost vizuálních formátů nutričních etiket napříč značkami a zeměmi, která může vést k~nižší úspěšnosti rozpoznání u~atypicky řešených tabulek.

\section{Implementace klíčových funkcí}
\label{sec:klicoveFunkce}

Tato sekce popisuje implementaci hlavních funkčních celků aplikace, které doplňují vstupní modality popsané v~sekci~\ref{sec:vstupniModality} o~přehledové, analytické, motivační a konfigurační funkce.
Každá podsekce odkazuje na příslušné funkční požadavky z~kapitoly~\ref{ch:navrh} a popisuje konkrétní implementační detaily včetně příslušných obrazovek aplikace.

\subsection{Denní přehled a správa nutričních cílů}
\label{subsec:dennPrehled}

Denní přehled (FR-02) je centrálním prvkem aplikace, který je realizován na hlavní obrazovce (viz obrázek~\ref{fig:screenDashboard}) prostřednictvím dvojice kontrolerů spravujících aktuálně vybraný den.
Kontrolery společně načítají odpovídající denní záznam včetně všech jídel a cvičení a propagují agregované nutriční hodnoty do reaktivních proměnných sledovaných uživatelským rozhraním.

Správa nutričních cílů (FR-03) je centralizována do dedikované služby, která vychází z~Mifflin-St~Jeorovy rovnice doplněné o~uživatelovu úroveň aktivity a cíl (hubnutí, udržení, nabírání).
Cíle jsou perzistovány prostřednictvím repozitáře pro denní záznamy a propagovány na aktuální i budoucí dny.
Konfigurace probíhá na obrazovce nutričních cílů v~profilové sekci (viz obrázek~\ref{fig:screenNutritionGoals}).

\begin{figure}[ht]
	\centering
	\screenshot[width=0.42\textwidth]{figures/screens/ProfileScreen_NutritionGoalsScreen.png}
	\caption{Nastavení nutričních cílů.}
	\label{fig:screenNutritionGoals}
\end{figure}

Doplňkovými funkcemi denního přehledu jsou systém přenosu nevyčerpaných kalorií do následujícího dne (rollover, omezený na nejvýše 500~kcal) a vizualizace plnění cílů v~kalendáři pomocí barevných kruhů znázorňujících podíly jednotlivých makroživin na cíli.
Kalendář na hlavní obrazovce je horizontálně scrollovatelný a umožňuje rychlou navigaci mezi dny bez ztráty kontextu denního přehledu.

\subsection{Správa jídel a ingrediencí}
\label{subsec:spravaJidel}

Správa jídel a ingrediencí (FR-01, FR-13, FR-15) je realizována na editační obrazovce jídla a navazujících obrazovkách pro editaci a detail jednotlivých ingrediencí.
Tyto obrazovky sdílí společnou logiku editace mezi rozpoznanými jídly z~pipeline umělé inteligence i čistě manuálně vytvořenými záznamy, čímž je zajištěna jednotnost uživatelského rozhraní.

Funkce \uv{Fix with AI} (FR-13) umožňuje z~editační obrazovky znovu vyvolat rozpoznání pomocí umělé inteligence s~doplněným kontextem, který uživatel mezitím přidal (například upřesnění typu pokrmu, korekce gramáže nebo poznámka k~jídlu).
Tok je realizován přes samostatnou obrazovku porovnání, která zobrazuje původní a nový výsledek vedle sebe a umožňuje uživateli zvolit, kterou verzi přijmout.

Pro funkci duplikace záznamu jídla na jiný den (FR-18) je implementován modální list s~kalendářovým výběrem dne.
Tento tok byl explicitně požadován respondenty uživatelské studie pro situace, kdy uživatel jí podobné jídlo opakovaně v~průběhu týdne.
Podpora jednotek (FR-15) zahrnuje gramy, 100-gramové porce, vlastní pojmenované jednotky a zobrazení zlomků (například 1/2 porce, 3/4 šálku).

\subsection{Sledování cvičení}
\label{subsec:sledovaniCviceni}

Modul cvičení naplňuje funkční požadavek FR-22 (částečně FR-23) a je realizován dedikovanou obrazovkou (viz obrázek~\ref{fig:screenExerciseLog}), která sjednocuje čtyři způsoby záznamu fyzické aktivity: rozpoznání pomocí umělé inteligence z~textového popisu, hlasový vstup, ruční zadání a výběr ze šablon dříve zaznamenaných cvičení.
Vlastní vstup probíhá na navazující obrazovce, která pro automatickou analýzu textových nebo hlasových popisů využívá centrální pipeline pro umělou inteligenci (viz sekce~\ref{sec:integraceAi}).

Detail jednotlivého záznamu cvičení je dostupný na samostatné obrazovce, doplněné o~modální list pro pokročilé akce (úprava, odstranění, převedení do šablon).
Šablony cvičení jsou modelovány samostatnou entitou v~datovém modelu (viz sekce~\ref{subsec:sablony}) a poskytují uživateli rychlý přístup k~opakovaně používaným aktivitám bez nutnosti opětovného volání umělé inteligence.

Energetický výdej spočítaný z~rozpoznaného nebo zadaného cvičení je propagován do hlavního denního přehledu, kde je zobrazen vedle příjmu energie pro splnění funkčního požadavku FR-22 (příjem versus výdej v~jednom pohledu).

\subsection{Sledování hmotnosti}
\label{subsec:sledovaniHmotnosti}

Sledování tělesné hmotnosti je realizováno samostatným podsystémem nezávislým na denním nutričním cyklu.
Záznam hmotnosti probíhá v~modálním listu (viz obrázek~\ref{fig:screenWeightSheet}), který umožňuje zadání aktuální hodnoty s~volitelnou fotografií dokumentující průběh redukce nebo nárůstu.
Historie záznamů a jejich úprava jsou dostupné na samostatných obrazovkách.

\begin{figure}[ht]
	\centering
	\screenshot[width=0.42\textwidth]{figures/screens/ProgressScreen_WeightBottomSheet.png}
	\caption{Záznam tělesné hmotnosti.}
	\label{fig:screenWeightSheet}
\end{figure}

Perzistenci a reaktivní aktualizace dat zajišťuje samostatný repozitář, který poskytuje vyšším vrstvám aplikace proud záznamů hmotnosti pro vizualizaci.
Karty s~aktuální hodnotou indexu BMI vypočteného podle rovnice~(\ref{eq:bmi}) a s~grafem vývoje hmotnosti za zvolené období jsou součástí obrazovky přehledů (viz obrázek~\ref{fig:screenProgress}).
Modul je propojen s~onboardingem aplikace, kde uživatel zadává počáteční a cílovou hmotnost spolu s~požadovanou rychlostí redukce nebo nárůstu.

\subsection{Oblíbené položky, šablony a historie}
\label{subsec:oblibene}

Rychlost zápisu opakovaných jídel a cvičení (FR-17, FR-18, FR-19), identifikovaná v~uživatelské studii jako klíčový faktor dlouhodobé udržitelnosti, je řešena třemi komplementárními mechanismy.
Oblíbené položky jsou modelovány samostatným příznakem na entitách jídla, ingredience a cvičení a takto označené položky jsou prezentovány v~záložce \uv{Oblíbené} na obrazovce výběru jídla (viz obrázek~\ref{fig:screenMealLog}).
Systém šablon (sekce~\ref{subsec:sablony}) doplňuje oblíbené o~automaticky udržovanou kanonickou verzi opakovaně používaných položek s~metrikou počtu použití pro řazení podle frekvence.
Našeptávání podle historie (FR-19) je realizováno fulltextovým vyhledáváním s~debounce.
Duplikace záznamu na jiný den (FR-18) využívá modální list popsaný v~sekci~\ref{subsec:spravaJidel}.

\subsection{Dietní omezení a vizualizace porušení}
\label{subsec:dietniOmezeni}

Dietní kontext (FR-20) eviduje uživatelská relace ve dvou typech nastavení: předdefinovaný typ diety (vegetariánská, veganská, bezlepková apod.) a vlastní omezení ve volném textovém seznamu.
Při každé analýze jídla je tento kontext zahrnut do promptu pro umělou inteligenci společně s~obrazovým nebo textovým vstupem.
Samotné vyhodnocení, zda rozpoznané jídlo porušuje nastavená uživatelská omezení, je svěřeno modelu, který případnou kolizi označí přímo ve své strukturované odpovědi.
Vizualizace porušení v~kalendáři (FR-21) je realizována kartou na obrazovce přehledů s~měsíční mřížkou a barevně odlišenými dny, kliknutím lze přejít na detailní pohled.

\subsection{Dotazy v~přirozeném jazyce}
\label{subsec:askAi}

Dotazy v~přirozeném jazyce (FR-27) jsou realizovány prostřednictvím dedikované obrazovky a kontroleru (viz obrázek~\ref{fig:screenAskAi}), které poskytují uživateli analytické rozhraní nad jeho historickými nutričními daty.

Implementace využívá dvouprůchodový systém pro snížení nákladů a zlepšení kvality odpovědí.
V~prvním průchodu volá pipeline menší instanci modelu s~úkolem identifikovat časový rozsah dotazu a typ analytické operace (například \uv{minulý týden}, \uv{leden 2026}, \uv{období bez logování}).
Ve druhém průchodu je hlavnímu modelu předán pouze relevantní výřez dat z~lokální databáze prostřednictvím repozitářové vrstvy, čímž se snižuje velikost kontextu a zlepšuje přesnost odpovědi.
Tento návrh umožňuje uživateli pokládat různorodé dotazy od trendů a srovnání období přes detekci vzorců stravování až po doporučení pro úpravu jídelníčku.

\subsection{Notifikace a motivační souhrny}
\label{subsec:notifikace}

Připomínky stravování (FR-28, FR-29) jsou realizovány dedikovanou službou postavenou nad knihovnou \texttt{flutter\_local\_notifications}, která plánuje notifikace s~ohledem na aktuální časové pásmo zařízení.
Aplikace nabízí pět typů stravovacích připomínek (snídaně, svačina, oběd, odpolední svačina, večeře) doplněných o~volitelné závěrečné shrnutí dne.
Jejich časy a aktivaci uživatel konfiguruje na samostatné obrazovce nastavení (viz obrázek~\ref{fig:screenReminders}).

Doplňkovou funkcí je generování motivačních souhrnů (FR-28), které periodicky analyzují data uživatele (denně, týdně, měsíčně) a prostřednictvím modelu umělé inteligence vytvářejí krátké personalizované shrnutí jako pozitivní zpětnou vazbu pro dlouhodobé cíle.
Souhrny jsou zobrazovány v~modálním listu (viz obrázek~\ref{fig:screenMotivationalSummary}).

\begin{figure}[ht]
	\centering
	\begin{minipage}[t]{0.31\textwidth}
		\centering
		\screenshot[width=\textwidth]{figures/screens/ProfileScreen_TrackingRemindersScreen.png}
		\\(a)~Přehled připomínek
	\end{minipage}\hfill
	\begin{minipage}[t]{0.31\textwidth}
		\centering
		\screenshot[width=\textwidth]{figures/screens/ProfileScreen_TrackingReminders_BottomSheet.png}
		\\(b)~Nastavení času
	\end{minipage}\hfill
	\begin{minipage}[t]{0.31\textwidth}
		\centering
		\screenshot[width=\textwidth]{figures/screens/MotivationalSummaryScreen_BottomSheet.png}
		\\(c)~Motivační souhrn
	\end{minipage}
	\caption{Notifikace a motivační souhrny.}
	\label{fig:screenReminders}
	\label{fig:screenMotivationalSummary}
\end{figure}

\subsection{Export dat a sdílení}
\label{subsec:exportDat}

Export dat (FR-25) podporuje dva výstupní formáty: tabulkový CSV pro další zpracování v~externích nástrojích a formátovaný PDF report pro archivaci nebo sdílení s~odborníkem na výživu.
Tok exportu začíná na úvodní obrazovce (viz obrázek~\ref{fig:screenExport}), kde uživatel zvolí formát, časový rozsah a obsah, a pokračuje na navazujících obrazovkách pro detailní konfiguraci PDF výstupu.

\begin{figure}[ht]
	\centering
	\begin{minipage}[t]{0.42\textwidth}
		\centering
		\screenshot[width=\textwidth]{figures/screens/ExportSummaryReportScreen.png}
		\\(a)~Konfigurace exportu
	\end{minipage}\hfill
	\begin{minipage}[t]{0.42\textwidth}
		\centering
		\screenshot[width=\textwidth]{figures/screens/ExportSummaryReportScreenExport.png}
		\\(b)~Vygenerovaný export
	\end{minipage}
	\caption{Tok exportu dat.}
	\label{fig:screenExport}
\end{figure}

Sdílení individuálních záznamů jídla je realizováno samostatnou službou nad knihovnou \texttt{share\_plus}, která prostřednictvím systémového dialogu distribuuje snímek jídla nebo strukturovaný textový záznam.

\subsection{Integrace se zdravotními platformami}
\label{subsec:zdravotniIntegrace}

Integrace s~Apple Health (iOS) a Health Connect (Android) je realizována dedikovanou službou nad knihovnou \texttt{health}, která poskytuje jednotné API pro čtení a zápis zdravotních dat napříč oběma platformami.
Konfigurace integrace probíhá na samostatné obrazovce v~profilové sekci (viz obrázek~\ref{fig:screenHealthIntegration}).

\begin{figure}[ht]
	\centering
	\screenshot[width=0.42\textwidth]{figures/screens/ProfileScreen_SyncAppleHealthScreen.png}
	\caption{Integrace se zdravotními platformami.}
	\label{fig:screenHealthIntegration}
\end{figure}

Načtené údaje o~spálených kaloriích jsou promítány do denního přehledu.
Samostatný atribut zdroje odlišuje manuálně zadaná cvičení od importovaných.
Uživatel má prostřednictvím samostatných přepínačů kontrolu nad tím, zda a jak se importovaný energetický výdej promítá do denního kalorického limitu a zda se nevyčerpané kalorie přenášejí do následujícího dne.
Aktuální implementace neobsahuje granulární multiplikátory pro různé typy aktivit (FR-23 částečně, viz sekce~\ref{sec:omezeniImplementace}).

\subsection{Lokalizace, onboarding a profil uživatele}
\label{subsec:lokalizaceOnboarding}

Lokalizace aplikace (čeština, angličtina) je realizována knihovnou \texttt{easy\_localization} s~automaticky generovanými klíči a JSON překladovými soubory v~adresáři s~assets.
Onboardingový tok (FR-04, částečně FR-09 a FR-30) zahrnuje jedenáct obrazovek včetně volitelné obrazovky \uv{Custom Diet} pro vyplnění vlastních dietních preferencí a provází nového uživatele od zadání základních údajů přes volbu cíle a výpočet doporučených kalorií až po uložení prvního denního záznamu.
Profilová obrazovka osobních údajů a obrazovka předvoleb aplikace (viz obrázek~\ref{fig:screenPersonalDetailsAndPrefs}) slouží jako primární rozhraní pro pozdější úpravu profilu a doplňková nastavení.

\begin{figure}[ht]
	\centering
	\begin{minipage}[t]{0.42\textwidth}
		\centering
		\screenshot[width=\textwidth]{figures/screens/ProfileScreen_PersonalDetailsScreen.png}
		\\(a)~Osobní údaje
	\end{minipage}\hfill
	\begin{minipage}[t]{0.42\textwidth}
		\centering
		\screenshot[width=\textwidth]{figures/screens/ProfileScreen_Preferences.png}
		\\(b)~Předvolby aplikace
	\end{minipage}
	\caption{Profilové obrazovky aplikace.}
	\label{fig:screenPersonalDetailsAndPrefs}
\end{figure}

Atributy uživatelské relace jsou perzistovány prostřednictvím dedikované služby nad klíč-hodnota úložištěm.
FR-30 (skrytí nebo zobrazení pokročilých funkcí) je implementován individuálními přepínači v~předvolbách (přebírání spálených kalorií, přenos nevyčerpaných kalorií, automatická úprava makroživin), bez explicitního režimu \emph{basic} versus \emph{advanced}.

\subsection{Další implementované funkce}
\label{subsec:dalsiFunkce}

Vedle hlavních funkcí pokrývajících funkční požadavky obsahuje aplikace několik doplňkových mechanismů, které vznikly v~průběhu implementace na základě praktických potřeb a zpětné vazby z~průběžného testování.

Sledování série po sobě jdoucích dnů se záznamy je realizováno samostatným podsystémem, který udržuje aktuální a nejdelší dosaženou sérii a zobrazuje je v~motivačním dialogu při dosažení významných milníků.
Tato funkce slouží jako pozitivní zpětná vazba v~návaznosti na zjištění z~uživatelské studie, kde respondenti opakovaně zmiňovali význam motivačních prvků pro dlouhodobou udržitelnost.

Aplikace dále podporuje widgety na domovské obrazovce zařízení prostřednictvím dedikované synchronizační vrstvy a směrovače akcí, které umožňují uživateli vidět denní souhrn kalorií a otevřít konkrétní vstupní modalitu bez nutnosti spouštět aplikaci.
Vizualizace plnění cílů v~kalendáři využívá vlastní renderování kruhových indikátorů pro každý den, kde tloušťka jednotlivých kruhů odráží podíl jednotlivých makroživin na cíli.

Pro zlepšování kvality rozpoznávání pomocí umělé inteligence je k~dispozici funkce hlášení nesprávně rozpoznaného jídla, která sbírá uživatelskou zpětnou vazbu pro budoucí ladění promptů.
Funkce automatické úpravy makroživin proporcionálně přepočítává cíle pro bílkoviny, sacharidy a tuky při změně kalorického cíle, čímž odstiňuje uživatele od ručního dopočítávání podílů.

Z~hlediska funkčního požadavku FR-05 (kontrola nad daty) aplikace umožňuje mazání jednotlivých záznamů jídel a cvičení, neobsahuje však úplné mazání uživatelského profilu.
Funkční požadavek FR-26 (offline tolerance) je naplněn principem \emph{local-first}, kdy jsou veškerá uživatelská data primárně ukládána do lokální SQLite databáze (viz sekce~\ref{sec:datovyModel}), což zajišťuje plnou funkčnost aplikace bez internetového připojení s~výjimkou modalit závislých na cloudovém modelu umělé inteligence.

\section{Přehled implementovaných funkčních požadavků}
\label{sec:prehledFR}

Pro shrnutí míry pokrytí návrhové specifikace je v~tabulce~\ref{tab:fr-status} uveden přehled stavu implementace všech třiceti funkčních požadavků definovaných v~sekci~\ref{sec:pozadavky}.
Každý požadavek je klasifikován do jedné ze tří kategorií, kterými jsou implementováno, částečně implementováno a neimplementováno, doplněných o~stručné zdůvodnění a odkaz na příslušné komponenty aplikace.
Tato klasifikace tvoří přímý podklad pro vyhodnocení splnění návrhových cílů v~kapitole~\ref{ch:zaver} a pro identifikaci oblastí, jejichž ověření je předmětem uživatelského testování v~kapitole~\ref{ch:testovani}.

\begin{table}[ht]
	\centering
	\caption{Stav implementace funkčních požadavků.}
	\label{tab:fr-status}
	\small
	\begin{tabular}{p{1.1cm}p{5.5cm}p{2.6cm}p{5cm}}
		\toprule
		\textbf{FR} & \textbf{Název} & \textbf{Stav} & \textbf{Poznámka} \\
		\midrule
		FR-01 & Ruční přidání záznamu jídla bez umělé inteligence & Implementováno & Editační obrazovka jídla, manuální vstup ingrediencí \\
		FR-02 & Denní přehled kalorií a makroživin & Implementováno & Hlavní obrazovka s~kontrolery denního záznamu \\
		FR-03 & Nastavení cílových hodnot & Implementováno & Služba nutričních cílů s~obrazovkou nastavení \\
		FR-04 & Správa profilu uživatele & Implementováno & Profilové obrazovky, služba uživatelské relace \\
		FR-05 & Kontrola nad daty a mazání & Částečně implementováno & Mazání jídel a cvičení funkční, mazání účtu chybí \\
		FR-06 & Fotografie jako vstup pro rozpoznání & Implementováno & Obrazovka kamery, pipeline umělé inteligence \\
		FR-07 & Návrh položek a porcí umělou inteligencí & Implementováno & \texttt{AiPipelineService}, strukturovaný JSON prompt \\
		FR-08 & Indikace nejistoty výstupu & Implementováno & Barevný badge (zelená $\geq$ 75~\%, žlutá $\geq$ 50~\%, červená $<$ 50~\%) \\
		FR-09 & Vysvětlení limitů umělé inteligence & Implementováno & Pětistránková onboardingová sekvence s~tipy \\
		FR-10 & Rozlišení limitu umělé inteligence a technické chyby & Částečně implementováno & Čárový kód: 6 typů chyb, AI pipeline: generické selhání \\
		FR-11 & Textový fallback po selhání rozpoznání & Implementováno & Textový popis $\to$ pipeline umělé inteligence \\
		FR-12 & Import fotografie z~galerie & Implementováno & Integrace knihovny \texttt{image\_picker} \\
		FR-13 & Opakování rozpoznání z~editace & Implementováno & \uv{Fix with AI} a samostatná obrazovka porovnání verzí \\
		FR-14 & Záznam bez fotografie (hlas, text) & Implementováno & Hlasový vstup, textový popis, manuální vstup \\
		FR-15 & Jednotky množství (gramy, kusy) & Implementováno & 1~g, 100~g, vlastní jednotky, zlomky \\
		FR-16 & Čtečka čárových kódů & Implementováno & Knihovna \texttt{mobile\_scanner}, klient Open Food Facts \\
		FR-17 & Oblíbené položky & Implementováno & Příznak na jídlech, ingrediencích a cvičeních \\
		FR-18 & Duplikace předchozího záznamu & Implementováno & Modální list s~výběrem data \\
		FR-19 & Našeptávání podle historie & Částečně implementováno & Fulltextové vyhledávání s~debounce, bez klasického dropdown \\
		FR-20 & Dietní omezení a intolerance & Implementováno & Předdefinovaný typ diety + uživatelská omezení \\
		FR-21 & Porušení dietních omezení v~kalendáři & Implementováno & Karta s~měsíční mřížkou na obrazovce přehledů \\
		FR-22 & Příjem versus výdej v~jednom pohledu & Implementováno & Hlavní obrazovka s~kalorickým příjmem a výdajem \\
		FR-23 & Nastavení integrace výdeje & Částečně implementováno & Přepínače pro výdej a rollover, bez granulárních multiplikátorů \\
		FR-24 & Týdenní a měsíční přehledy & Implementováno & Obrazovka přehledů s~kartou týdenní energie \\
		FR-25 & Export dat (CSV, PDF) & Implementováno & Dedikovaná služba exportu s~výběrem období \\
		FR-26 & Offline tolerance & Implementováno & Lokální SQLite, umělá inteligence vyžaduje připojení \\
		FR-27 & Dotazy v~přirozeném jazyce & Implementováno & Obrazovka Ask AI, dvouprůchodový systém \\
		FR-28 & Měsíční motivační souhrn & Implementováno & Periodicky generované shrnutí (denně, týdně, měsíčně) \\
		FR-29 & Nastavitelné notifikace & Implementováno & Pět typů připomínek s~knihovnou \texttt{flutter\_local\_notifications} \\
		FR-30 & Skrytí a zobrazení pokročilých funkcí & Částečně implementováno & Individuální přepínače, bez režimu basic versus advanced \\
		\bottomrule
	\end{tabular}
\end{table}

\section{Omezení implementace}
\label{sec:omezeniImplementace}

Aplikace Foody představuje funkční prototyp ověřující navržený koncept a v~několika oblastech nedosahuje plného potenciálu návrhu z~kapitoly~\ref{ch:navrh}.
Tato sekce shrnuje hlavní omezení výsledné implementace, která je třeba zohlednit při interpretaci výsledků uživatelského testování v~kapitole~\ref{ch:testovani} a při formulaci doporučení pro další rozvoj v~kapitole~\ref{ch:zaver}.

Hlavní omezení vyplývá z~povahy rozpoznávání pomocí umělé inteligence.
Multimodální jazykový model dosahuje vyšší přesnosti u~jednoznačně viditelných pokrmů, zatímco u~složených jídel s~překrývajícími se ingrediencemi a u~regionálně specifických receptur klesá spolehlivost odhadu.
Kvalita rozpoznání je rovněž závislá na kvalitě vstupního obrazu, zejména na osvětlení a úhlu pohledu, což aplikace adresuje pětistránkovou onboardingovou instruktáží (FR-09), nemůže však zcela kompenzovat nepříznivé pořizovací podmínky.
Latence rozpoznávání se v~kontrolovaných podmínkách pohybuje hluboko pod stanoveným limitem 20~vteřin (NFR-02), avšak energetický odhad zůstává aproximací s~nejistotou, kterou aplikace transparentně signalizuje prostřednictvím tříúrovňového indikátoru spolehlivosti.

Architektura aplikace závisí na třech externích službách, jejichž dostupnost ovlivňuje uživatelskou zkušenost.
Komunikace s~poskytovatelem modelu umělé inteligence vyžaduje aktivní internetové připojení, takže funkce rozpoznávání jídel z~fotografie, hlasového a textového popisu, analýzy cvičení a dotazů v~přirozeném jazyce nejsou v~režimu offline dostupné, zatímco zbývající funkcionality jsou plně funkční lokálně (FR-26).
Vyhledávání podle čárového kódu se opírá o~veřejnou databázi Open Food Facts, jejíž pokrytí lokálních českých produktů není úplné, což bylo identifikováno jako omezení i v~uživatelském testování (sekce~\ref{subsec:uiProblemyTesting}).
Provoz multimodálního modelu představuje nezanedbatelný náklad na volání API, jehož výše je úměrná frekvenci použití funkcí postavených na umělé inteligenci a v~současné implementaci není limitován kvótami na úrovni uživatele.

Pět funkčních požadavků (FR-05, FR-10, FR-19, FR-23 a FR-30) bylo implementováno pouze částečně.
Konkrétní povaha omezení jednotlivých požadavků je uvedena ve sloupci \emph{Poznámka} tabulky~\ref{tab:fr-status} v~sekci~\ref{sec:prehledFR}.

Z~technických omezení je v~současné implementaci skenování nutričních etiket pokryto pouze pro základní makroživiny, zatímco rozšířené atributy jako vláknina, sůl a vybrané mikronutrienty zůstávají námětem pro budoucí iteraci.
Analýza cvičení a výpočet nutričních cílů jsou v~aktuální verzi vázány přímo na poskytovatele OpenAI a nepřepínají se spolu s~hlavním poskytovatelem rozpoznávání jídel.
Ověření kompatibility Gemini API pro tyto scénáře nebylo v~rámci práce dokončeno.
Aplikace dále neobsahuje automatický fallback mezi poskytovateli při selhání jednoho z~nich a nedisponuje cloudovou synchronizací ani vzdálenou zálohou dat, což odpovídá zvolenému architektonickému přístupu bez vlastního backendu.

Posledním okruhem omezení je rozsah uživatelského testování.
Implementace byla ověřena na malém vzorku tří participantů v~kontrolovaném prostředí, jak je popsáno v~kapitole~\ref{ch:testovani}, což odpovídá rozsahu kvalitativního testování v~rámci diplomové práce, ale neumožňuje zobecnění výsledků na širší populaci.
Dlouhodobá retence uživatelů, přesnost odhadů v~reálném používání a chování aplikace v~produkčním prostředí na různých zařízeních a operačních systémech zůstávají předmětem dalšího ověření, jež je formulováno jako doporučení pro budoucí validační studii v~sekci~\ref{sec:doporuceni}.
